\section{The sharp local estimate at standard lattice places}

% BEGIN P76V1
\paragraph{P76V1. The different bounds the mixed coefficient on the marked directions.}
Work on the original arithmetic host $E=D.\mathrm{host}.E$ and on
the test and integral chart of R76V1. Put $\nu=\mathrm{ord}_{V'}/e'$
with $e'=be$, let $\ell$ be the residue characteristic, and put
\[
 c_\ell=\nu(\ell)\frac{\operatorname{ord}_\ell(|D_E|)}
                             {[E:\mathbb Q]}.
\]
There is a canonical gauge for which the mixed coefficient
$B'(u,v)$ has $\nu(B'(u,v))\ge-c_\ell$ for every two of the
integral eigenvectors in the given frame. We allow all embeddings
of the coefficient field in the original valuation tests.

To see the arithmetic input, choose a $\mathbb Z$-basis $(x_i)$
of $\mathcal O_E$ and its trace-dual basis $(y_i)$. At a prime
above $\ell$, let $e_0,f_0,t$ be the ramification index, residue
degree and different exponent, and let $g$ be the number of primes
above $\ell$. The host is Galois, so these three numbers agree at
all these primes. Taking the norm of the different and using the
degree formula gives $\operatorname{ord}_\ell(|D_E|)=gf_0t$ and
$[E:\mathbb Q]=ge_0f_0$. Each $y_i$ belongs to the inverse
different, hence $\nu(\sigma(y_i))\ge-(t/e_0)\nu(\ell)=-c_\ell$
for every embedded character $\sigma$.

Start with the mixed coefficient $B$ of any canonical gauge.
For integral tangent vectors $u,v$ the symmetry defect
$B(v,u)-B(u,v)$ and all adjoint defects
$B(A(x_i)u,v)-B(u,A(\bar x_i)v)$ are integral, by the
mixed-coefficient identities of R76V3. Decompose the tangent
space over the chart field into its $E$-eigenspaces. Indeed the
coefficient field contains all the embeddings and $E$ is
separable. Write $B_*$ for the bilinear form which equals $B$
on $(\sigma,\bar\sigma)$ blocks and is zero on the other blocks.
It is balanced for the adjoint action. For eigenvectors of types
$\sigma,\tau$ with $\sigma\ne\tau\circ\mathrm{bar}$, trace
duality gives
\[
 1=\sum_i\sigma(y_i)
              \bigl(\sigma(x_i)-(\tau\circ\mathrm{bar})(x_i)\bigr).
\]
Indeed the two embedding matrices of trace-dual bases are inverse
transposes. Multiplying by $B(u,v)$ turns each parenthesized
factor into an integral adjoint defect. Thus $B(u,v)$ itself has
valuation at least $-c_\ell$ on these blocks.

Define $W=A(w_0)$ and
$C(u,v)= B_*(Wu,v)+B_*(Wv,u)$. This form is symmetric and
balanced. On complementary blocks the two coefficients of
$B(u,v)$ and $B(v,u)$ are $\sigma(w_0)$ and
$\sigma(\bar w_0)$, whose sum is one. Thus $B-C$ on such an
integral pair is integral by the symmetry defect. On all other
blocks $C$ vanishes. Subtract $C$ in the two formal logarithm
coordinates from the original gauge. Biadditivity, symmetry and
balance preserve the gauge equations as in G76J1. Its mixed
coefficient $B'=B-C$ has the asserted bound, including in
residue characteristic two.
% END P76V1

% BEGIN P76V2
\paragraph{P76V2. A loss can occur only on disjoint pairs of derivatives.}
For the gauge $h$ of P76V1 set $q(x)=h(x,[w_0]x)$. For any
weight-$k$ integral column germ $f$, a jet of order $j$ in
the prescribed integral eigendirections, computed from
$f\exp(kq)$, has valuation at least $-\lfloor j/2\rfloor c_\ell$.
Indeed $q(0)=0$ and its first derivatives vanish at zero. For
two such directions the unscaled integral identity proved in
R76V3 reads
\[
 D_uD_vq - B'(u,Wv)-B'(v,Wu)\ \in V'[[x]].
\]
Here the two terms displayed after the derivative are constants.
Each has valuation at least $-c_\ell$, since $W$ acts on integral
eigendirections by algebraic integers. Every further integral
invariant derivative annihilates those constants and preserves
integral series. Hence a derivative of $q$ at zero of order at
least three is integral.

Use labelled derivatives, allowing repeated directions. By the
product rule, choose a subset to differentiate $f$ and partition
its complement into nonempty blocks. The corresponding term at
zero is the chosen derivative of $f$ times, for each block, $k$
times that block's derivative of $q$. This follows by adjoining
one derivative: it either hits $f$, starts a singleton on the
exponential, or differentiates one existing block. There are no
factorial denominators. Only blocks of size two can have negative
valuation, and they are disjoint. This proves the jet estimate.
For $P$ integral at the test, multiply the column jets of a
monomial and then sum with its coefficients. Every nonzero
permitted $P$ has total degree $d=d_w$ by its row degrees, and
the sum of the column orders is $d$. Harmonic gauge invariance
therefore gives, for the original connection and reached columns,
\[
 \nu(\mathrm{hv}(F,P))\ge-\lfloor d/2\rfloor c_\ell
                 \ge -B_{\mathrm{ICL}},
 \qquad B_{\mathrm{ICL}}=\tfrac{d_w}{2}
          \sum_{p\mid\mathrm{condN}}\nu(p)\,\mathrm{vdH}(p).
\]
The Galois host has every discriminant prime in its conductor,
and $\mathrm{vdH}(p)=\operatorname{ord}_p(|D_E|)/[E:\mathbb Q]$.
Only $p=\ell$ can contribute. Zero polynomials and zero values
satisfy the inequality with the usual infinite valuation.
% END P76V2

% BEGIN P76V3
\paragraph{P76V3. Cancellation of the primary multiplicity and the auxiliary level.}
Let $S=\mathcal O_E^n$. Assume that
$\Lambda\otimes\mathbb Z_\ell=S\otimes\mathbb Z_\ell$ and that
$M$ is squarefree. For every original $q$, (ICL.14)--(ICL.15)
hold with their prescribed $\alpha$, for arbitrary $\ell$-part of
$N$ and arbitrary ramification of the valuation test.

The equality of local lattices implies that the coefficients of
the generators $\mathrm{letB}(\sigma,i,v)$ of $\mathrm{symO}$
are integral at the test. Indeed for $v\in\Lambda$ a rational
integer prime to $\ell$ clears its coordinates in $S$. All
embedded coefficients of $P\in\mathrm{harmL}$ are consequently
integral. Write $N=m\ell^s$ with $m$ prime to $\ell$. Split the
finite group of enumerated isotropic tuples as $Z_\ell\times Z'$,
with $Z'$ prime to $\ell$. The local lattice gives
$|Z_\ell|=\ell^{\operatorname{ord}_\ell(\mathrm{mIdx})}$:
the full order is $[\Lambda:NS]^r$, and the quotient by
$[\Lambda\cap S:NS]^r=\mathrm{mIdx}$ is
$[\Lambda:\Lambda\cap S]^r$, prime to $\ell$. Every coordinate
of a tuple in $Z_\ell$ is killed by $N$. In fact there is an
integer $a$ prime to $\ell$ with $a\Lambda\subset S$. Apply
the defining matrix test of $\mathrm{InZh}$ with $aI$ to see
that $aN$ kills each coordinate, and remove the prime-to-$\ell$
factor on the primary subgroup.

The symmetric translation at such a coordinate fixes the
normalized $[N]$-pullback of the generator: its deck action is
the pullback of the identity translation, and odd symmetry
selects that lift. The lift for a product of primary and
prime-to-$\ell$ points is the product of their lifts. Their
commutator has order dividing both orders, and their commuting
product is the unique symmetric lift. The last translation, by
$t_q$, is applied to the same section after this equality.
Every $Z'$ term is therefore repeated exactly $|Z_\ell|$ times.
This cancels the whole $\ell$-part of $\mathrm{mIdx}$.

Put $\alpha=\ell$ if $\ell\mid M$, and $\alpha=1$ otherwise.
Set $n_0=m\alpha$ and $b_0=N/n_0$. These are odd integers
because $M\mid N$ and $N$ is odd. For each surviving column
let $z'$ be its prime-to-$\ell$ point. Both $[b_0]z'$ and
$[b_0]t_q$ are killed by $n_0^2$. Away from $\ell$ this follows
from the original $N^2$-torsion. At $\ell\mid M$, the point
$Nt_q$ is $M$-torsion, so the primary order of $t_q$ divides
$\ell^{s+1}$; $b_0=\ell^{s-1}$ reduces it to at most $\ell^2$.
At $\ell\nmid M$, half-ideal killing gives $M^j t_q=0$ for
some $j$, so $t_q$ has no $\ell$-part.

At the lower power $n_0^2$ form the normalized $[n_0]$-pullback
followed, in order, by the symmetric translations at $[b_0]z'$
and $[b_0]t_q$. Denote the resulting column by $H$. These are
reachability operations of R76V2, so $H$ has an integral germ.
Its normalized $[b_0]$-pullback is exactly the original column.
Multiplication pullbacks compose by their common rigidification.
Translation commutes with pullback because $[b_0]$ intertwines
the two translations and pulls their odd symmetric lifts to the
same lifts in the pullback comparison. Equivalently, apply this
identity first to $z'$, then to $t_q$, on the common rigidified
line. The harmonic chain rule multiplies the value by $b_0^d$:
in formal logarithm coordinates $[b_0]$ multiplies each tangent
by $b_0$, its biextension pullback multiplies the tensor power
by $b_0^2$, and the transported canonical gauge gives the same
harmonic evaluation. Thus
\[
 \alpha^{d_w}X_q
  = \frac{|Z_\ell|}{\mathrm{mIdx}}\,m^{-d}
        \sum_{z'\in Z'}\mathrm{hv}(H_{q,z'},P).
\]
For $P\ne0$ this uses $\alpha^{d_w}b_0^d/N^d=m^{-d}$.
The prefactor and $m$ are units at the test. P76V2 and the
valuation inequality for sums prove the claim with the
normalization $\nu=\mathrm{ord}_{V'}/(be)$.

This includes every $\ell\mid M$ by the stipulated
$\mathrm{TrivAt}(M)$, and every $\ell\nmid N$ because
$NS\subset\Lambda\subset N^{-1}S$ by unimodularity. Thus the
remaining case of the sharp estimate consists of auxiliary
primes $\ell\mid N$, $\ell\nmid M$ where the local lattice differs
from $S\otimes\mathbb Z_\ell$.
% END P76V3
